\chapter{Reality measure, Pfaffian и статус orbit half-trace}

Cross-tome audit оставил условным правило
\begin{equation}
 \Tr_{\mathrm{orb}}=\frac12\Tr_{H_8}.
\end{equation}
Для закрытия нужно различить fermionic functional measure и bosonic
spectral action: они не обязаны использовать один и тот же trace
prescription.

\section{KO6 reality doubling}

Пусть
\begin{equation}
 H_8=H_4\oplus H_4^c,
\end{equation}
а antiunitary \(J\) обменивает две копии. Для \(J\)-compatible observable
\begin{equation}
 Q_8=Q_4\oplus JQ_4J^{-1}
\end{equation}
и вещественного spectral invariant получаем
\begin{equation}
 \Tr_{H_8}Q_8=2\Tr_{H_4}Q_4.
\end{equation}
Это algebraic doubling identity, но само по себе оно не определяет,
какой trace должен входить в каждый functional.

\section{Fermionic Pfaffian measure}

Для real even spectral data fermionic quadratic action имеет
антисимметричную форму
\begin{equation}
 S_F[\psi]
 =
 \frac12\langle J\psi,\mathcal D\psi\rangle.
\end{equation}
Grassmann integration даёт
\begin{equation}
 Z_F=\operatorname{Pf}(\mathcal C\mathcal D),
\end{equation}
и
\begin{equation}
 \log Z_F
 =
 \frac12\Tr\log(\mathcal C\mathcal D)
\end{equation}
на doubled complex representation. Поэтому particle--antiparticle
reality orbit считается один раз.

\begin{theorem}[Pfaffian half-count pass]
Для fermionic determinant, state multiplicities и one-loop fermion
contribution коэффициент \(1/2\) выводится из Pfaffian measure. Две
physical Dirac-пары finite ledger не должны удваиваться из-за формального
KO6 completion.
\end{theorem}

\section{Bosonic spectral trace}

Классический bosonic spectral action имеет вид
\begin{equation}
 S_B=\Tr_{H_8}f(\mathcal D^2/\Lambda^2)
\end{equation}
и использует полный complex trace. Fermionic Pfaffian не заменяет его
автоматически на half-trace.

Для doubled finite block:
\begin{equation}
 G_{\mathrm{full}}=2G_{\mathrm{orb}},
 \qquad
 H_{\mathrm{full}}=2H_{\mathrm{orb}}.
\end{equation}
Поэтому physical mass matrix invariant:
\begin{equation}
 (\kappa G_{\mathrm{full}})^{-1}H_{\mathrm{full}}
 =
 (\kappa G_{\mathrm{orb}})^{-1}H_{\mathrm{orb}}.
\end{equation}

Gauge normalization также invariant. Full trace даёт
\begin{equation}
 \Tr_{H_8}Q^2=4
\end{equation}
вместо orbit value \(2\), но scalar kinetic coefficient одновременно
удваивается с \(2\) до \(4\). Нормированный gauge coefficient остаётся
\begin{equation}
 \frac{(2/3)\cdot4}{4}=\frac23,
\end{equation}
следовательно,
\begin{equation}
 g^2=\frac38
\end{equation}
не меняется.

\begin{theorem}[Full-trace ratio invariance]
Переход от orbit-normalized bosonic notation к полному spectral trace
удваивает весь bosonic action глобально, но сохраняет \(\kappa\), scalar
masses, gauge coupling и vacuum orbit после canonical normalization.
\end{theorem}

\section{Исправленная trace dictionary}

Отныне различаем:
\begin{align}
 \Tr_{\mathrm{spec}}
 &:=\Tr_{H_8}
 &&\text{для bosonic spectral action},
 \\
 \Tr_{\mathrm{phys}}
 &:=\frac12\Tr_{H_8}
 &&\text{для Pfaffian fermion count},
 \\
 \Tr_{\mathrm{red}}
 &:=\Tr_{H_4}
 &&\text{как сокращённую вычислительную запись}.
\end{align}
Для \(J\)-invariant even polynomials
\begin{equation}
 \Tr_{\mathrm{red}}
 =\frac12\Tr_{\mathrm{spec}},
\end{equation}
но это identity representation, а не универсальный measure axiom.

Все прежние записи \(\Tr_{\mathrm{orb}}\) в normalized bosonic ratios
следует понимать как использование reduced representative с последующим
удалением общего doubling factor. В absolute bosonic action должен
использоваться \(\Tr_{\mathrm{spec}}\).

\section{Supertrace ledger}

Coleman--Weinberg supertrace считает physical propagating fields:
\begin{enumerate}
 \item три real scalar modes;
 \item три massive-vector polarizations;
 \item две physical Dirac pairs, посчитанные Pfaffian/reality measure.
\end{enumerate}
Следовательно,
\begin{equation}
 N_{\mathrm{fin+gauge}}
 =3\cdot4^2+3\cdot3^2-2\cdot4
 =67
\end{equation}
сохраняется. Формальное удвоение representation не создаёт новые
propagating particle species.

\section{Что понижается в статусе}

Role-graded norms, использовавшие absolute finite ranks
\begin{equation}
 \Tr_{\mathrm{orb}}P_H=23,
\end{equation}
не получают bosonic measure theorem из Pfaffian. Они уже исключены из
blind scorecard и сохраняются только как reduced-module construction
targets.

Аналогично абсолютная общая normalizing constant bosonic action зависит
от spectral moments и train scale. Но relative hidden-sector predictions
не зависят от выбора full или reduced trace.

\begin{nogocriterion}[No universal half-trace]
Запрещено использовать fermionic Pfaffian как доказательство
\(\frac12\Tr_{H_8}\) для независимого classical bosonic spectral action.
Half-count допустим для physical fermions; bosonic coefficients должны
считаться full trace либо явно нормироваться общим global factor.
\end{nogocriterion}

\begin{theorem}[Measure closure verdict]
Универсальный orbit half-trace отвергается, но основной hidden parent
action сохраняется. Fermionic half-count строго выводится из Pfaffian,
bosonic action использует full spectral trace, а все ранее выведенные
dimensionless ratios invariant относительно общего KO6 doubling.
\end{theorem}

\begin{protocol}[Следующий гейт]
Measure gap закрыт для внутренней hidden EFT. Следует проверить portal
menu при новой trace dictionary:
\begin{enumerate}
 \item scalar portal;
 \item abelian kinetic mixing;
 \item neutrino portal.
\end{enumerate}
Маршрут проходит только если symmetry или common algebra выбирает portal
и его normalization; произвольный continuous portal coefficient
сохраняет hidden sector phenomenologically disconnected.
\end{protocol}

Машинный аудит сохранён в
\path{s2t_v3_orbit_measure_pfaffian_results.json}.